% Laget automatisk av elevark/lag_elevark.py fra opplegg/fritt-fall-mobil.html. Endringer her blir overskrevet.
\documentclass[11pt]{article}
\usepackage{elevark}
\begin{document}
\elevarkhode{Fysikk 2 · Generell relativitetsteori}{Fritt fall og ekvivalensprinsippet med mobilen}{https://fysikklab.no/opplegg/fritt-fall-mobil.html}{Opplegget på nett}

Når mobilen ligger på bordet, viser akselerometeret 9,8~m/s². Når den faller fritt, viser det nesten null, selv om mobilen akselererer nedover med $g$. Mobilen kan ikke skille fritt fall fra å sveve i rommet. Det er ekvivalensprinsippet, og dere kan måle det i klasserommet.

\section*{Hensikt}

\begin{itemize}
\item Se at et akselerometer i fritt fall viser null, og forklare hvorfor.
\item Bestemme $g$ fra falltiden.
\item Knytte resultatet til ekvivalensprinsippet i generell relativitetsteori.
\end{itemize}

\section*{Utstyr}

\begin{itemize}
\item Mobil med appen phyphox (gratis) og eksperimentet «Akselerasjon med $g$».
\item En tykk pute, madrass eller dyne å slippe mobilen ned på. Slipp aldri mobilen på gulvet.
\item Målebånd. Til demonstrasjonen: en plastflaske med et par små hull nær bunnen, full av vann.
\end{itemize}

\section*{Teorien}

Et akselerometer måler kraften som skal til for å holde en liten masse inni sensoren på plass, delt på massen. På bordet må mobilen dytte massen oppover med en kraft som er like stor som tyngden, og sensoren viser 9,8~m/s². I fritt fall faller sensoren og massen like fort, og ingenting trenger å dytte. Sensoren viser 0. Inne i mobilen er det ingen forskjell på fritt fall og å sveve i et rom uten tyngdekraft. Ekvivalensprinsippet sier at ingen forsøk kan avgjøre forskjellen lokalt.

\section*{Demonstrasjon: vannflaska}

Hold flaska over en bøtte eller ute. Vannet spruter ut av hullene. Slipp flaska: Mens den faller, slutter vannet å renne. Vannet og flaska faller like fort, og vannet «veier» ingenting i forhold til flaska.

\section*{Framgangsmåte}

\begin{enumerate}
\item Start «Akselerasjon med $g$» og legg mobilen på bordet. Hva viser den? Hold den i hånda og løft den raskt opp og ned.
\item Hold mobilen flatt med skjermen opp 50~cm over puta. Start målingen og slipp. Se på grafen for den totale akselerasjonen: Den faller til nesten null mens mobilen faller, og får en topp når den treffer puta.
\item Les av hvor lenge akselerasjonen er nær null. Det er falltiden $t$. Mål fallhøyden fra undersiden av mobilen til toppen av puta.
\item Gjenta fra 70, 90 og 110~cm. Skriv inn i tabellen.
\item Ekstra: Kast mobilen forsiktig rett opp noen centimeter og ta den imot. Er akselerasjonen også null på vei opp?
\end{enumerate}

\section*{Oppgaver}

\begin{oppgaveliste}
\oppgave \fase{gjett}{Gjett}Hva viser akselerometeret mens mobilen faller: 9,8~m/s², 0 eller noe annet? Hva viser det når mobilen er kastet opp og er på vei oppover?
\svar
\oppgave \fase{test}{Mål}Gjør målingene, eller trykk på «Fyll inn eksempel». Hvilken verdi får dere for $g$?
\svar
\oppgave \fase{forklar}{Forklar}Mobilen akselererer med $g$ nedover, men akselerometeret viser null. Forklar hvordan det henger sammen.
\svar
\oppgave \fase{forklar}{Forklar}Hvorfor slutter vannet å renne ut av flaska mens den faller? Hva har det med astronautene i romstasjonen å gjøre?
\svar
\oppgave \fase{forklar}{Vurder}Mobilen viser sjelden helt null i fallet. Hva kan være grunnen? Hvordan påvirker det avlesningen av falltiden?
\svar
\oppgave \fase{utfordring}{Utfordring}Ekvivalensprinsippet sier at fritt fall i et tyngdefelt og å sveve i rommet ikke kan skilles fra hverandre. En astronaut i et lukket romskip i fritt fall rundt jorda har en svært nøyaktig akselerometer i hver ende av romskipet. Kan hun likevel finne ut at hun er nær jorda?
\svar
\end{oppgaveliste}

\end{document}
